% Compile with: latexmk -xelatex slides.tex
%
% The theme is built here rather than taken from a gallery, for two reasons. The
% Discrete Mathematics template puts frame titles in a full-width coloured bar,
% and this deck deliberately does not look like it. And metropolis, the obvious
% off-the-shelf alternative, is not in this TeX Live installation, so using it
% would force every member of the group to install a package before they could
% compile. Everything below ships with TeX Live.
%
% 19 frames for the 20-minute slot, following section 10.2 of
% KE_HOACH_TRIEN_KHAI.md.
\documentclass[11pt, aspectratio=169]{beamer}

% Shared preamble for report.tex and slides.tex.
% Compile both with XeLaTeX: latexmk -xelatex <file>.tex
%
% Visual identity, chosen to be distinct from the Discrete Mathematics template:
%   text     TeX Gyre Pagella (a Palatino), wider and rounder than Times
%   headings TeX Gyre Heros (a Helvetica), sans, to separate structure from prose
%   maths    TeX Gyre Pagella Math through unicode-math, so the digits inside
%            formulas match the digits in the surrounding sentence
% All three ship with TeX Live, so nothing extra has to be installed.

\usepackage{fontspec}
\usepackage{polyglossia}
\setmainlanguage{vietnamese}

\setmainfont{texgyrepagella}[
  Extension      = .otf,
  UprightFont    = *-regular,
  BoldFont       = *-bold,
  ItalicFont     = *-italic,
  BoldItalicFont = *-bolditalic,
]
\setsansfont{texgyreheros}[
  Extension      = .otf,
  UprightFont    = *-regular,
  BoldFont       = *-bold,
  ItalicFont     = *-italic,
  BoldItalicFont = *-bolditalic,
  Scale          = 0.92,
]
\setmonofont{texgyrecursor}[
  Extension      = .otf,
  UprightFont    = *-regular,
  BoldFont       = *-bold,
  Scale          = 0.88,
]

\usepackage{amsmath}
\usepackage{amssymb}
\usepackage{amsthm}
\usepackage{mathtools}
\usepackage{unicode-math}
\setmathfont{texgyrepagella-math.otf}

\usepackage{graphicx}
\usepackage{booktabs}
\usepackage{siunitx}
\usepackage{xcolor}

% Vietnamese decimal marker, so \num agrees with the decimals written by hand in
% the prose (1{,}39 rather than 1.39).
\sisetup{output-decimal-marker={,}}

% Slide copies come first on purpose: \resultgraphic must prefer them.
\graphicspath{{figures/}{../results/figures/slides/}{../results/figures/}}

% ---------------------------------------------------------------------------
% Palette. The three accents are the figure colours from src/figures.py, so a
% method named in the text can be tinted to match its curve.
% ---------------------------------------------------------------------------
\definecolor{ink}{HTML}{1A1A1A}
\definecolor{accent}{HTML}{0B5563}     % deep teal, the structural colour
\definecolor{rule}{HTML}{9AA5A8}
\definecolor{cGD}{HTML}{1F77B4}
\definecolor{cSGD}{HTML}{FF7F0E}
\definecolor{cAGD}{HTML}{2CA02C}
\definecolor{cNewton}{HTML}{D62728}

% ---------------------------------------------------------------------------
% Notation, kept identical to KE_HOACH_TRIEN_KHAI.md
% ---------------------------------------------------------------------------
\newcommand{\R}{\mathbb{R}}
\newcommand{\E}{\mathbb{E}}
\newcommand{\wstar}{w^{*}}
\newcommand{\fstar}{f^{*}}
\newcommand{\bigO}{\mathcal{O}}
\newcommand{\Batch}{\mathcal{B}}
\newcommand{\trans}{^{\top}}
\DeclareMathOperator*{\argmin}{arg\,min}
\DeclareMathOperator{\sign}{sign}
\DeclarePairedDelimiter{\norm}{\lVert}{\rVert}
\DeclarePairedDelimiter{\abs}{\lvert}{\rvert}

% ---------------------------------------------------------------------------
% Figures. A stem that has not been generated yet leaves a visible placeholder
% rather than a compilation error, so the report can be drafted while the
% experiments are still running.
% ---------------------------------------------------------------------------
\newcommand{\missingfigbox}[1]{%
  \fcolorbox{rule}{white}{\parbox[c][3.2cm][c]{0.86\linewidth}{\centering
    \ttfamily\small\detokenize{#1}.pdf\par\vspace{4pt}
    \rmfamily\small chưa được sinh ra từ notebook\par}}%
}

% The path is spelled out rather than left to \graphicspath, which lists the
% slide directory first and would otherwise hand the report the wide, reduced
% copies meant for Beamer.
%
% Placement is [tbp] with no 'h': a figure declared between two paragraphs of one
% argument would otherwise be typeset right there and cut the argument in half.
\newcommand{\resultfig}[3]{%
  \begin{figure}[tbp]
    \centering
    \IfFileExists{../results/figures/#1.pdf}
      {\includegraphics[width=0.82\linewidth]{../results/figures/#1.pdf}}
      {\missingfigbox{#1}}
    \caption{#2}
    \label{#3}
  \end{figure}%
}

% The mandatory pair, stacked rather than side by side. The source figures are
% 6 by 4 inches with 10 pt labels; at half the text width that becomes about 5 pt
% and the legends stop being readable, which defeats the purpose of printing two
% panels. Stacking costs most of a page per comparison and keeps both legible.
\newcommand{\panellabel}[1]{{\sffamily\small\color{accent}#1}}
\newcommand{\resultpair}[3]{%
  \begin{figure}[tbp]
    \centering
    \panellabel{(a) theo số vòng lặp}\par\vspace{2pt}
    \IfFileExists{../results/figures/#1_iters.pdf}
      {\includegraphics[width=0.8\linewidth]{../results/figures/#1_iters.pdf}}
      {\missingfigbox{#1_iters}}
    \par\vspace{10pt}
    \panellabel{(b) theo thời gian chạy}\par\vspace{2pt}
    \IfFileExists{../results/figures/#1_time.pdf}
      {\includegraphics[width=0.8\linewidth]{../results/figures/#1_time.pdf}}
      {\missingfigbox{#1_time}}
    \caption{#2}
    \label{#3}
  \end{figure}%
}

% Slides: figures are 6 by 4 inches, so a full-width copy is too tall for a frame
% that also carries text. Height is capped as a fraction of the text height and
% the aspect ratio kept; crowded frames pass a smaller fraction.
\newcommand{\resultgraphic}[2][0.72]{%
  \IfFileExists{../results/figures/slides/#2.pdf}
    {\includegraphics[width=\linewidth,height=#1\textheight,keepaspectratio]{../results/figures/slides/#2.pdf}}
    {\IfFileExists{../results/figures/#2.pdf}
      {\includegraphics[width=\linewidth,height=#1\textheight,keepaspectratio]{../results/figures/#2.pdf}}
      {\missingfigbox{#2}}}%
}


\usepackage{algorithm}
\usepackage{algpseudocode}
\floatname{algorithm}{Thuật toán}

% ---------------------------------------------------------------------------
% A flat theme: no shadows, no navigation, no sidebar. Structure is carried by
% one thin rule under the frame title and by the sans face, not by filled blocks.
% ---------------------------------------------------------------------------
\usetheme{default}
\usecolortheme{default}
\setbeamertemplate{navigation symbols}{}
\setbeamerfont{frametitle}{family=\sffamily, size=\large, series=\bfseries}
\setbeamerfont{framesubtitle}{family=\sffamily, size=\small}
\setbeamerfont{title}{family=\sffamily, size=\LARGE, series=\bfseries}
\setbeamerfont{author}{family=\sffamily, size=\small}
\setbeamerfont{date}{family=\sffamily, size=\small}
\setbeamercolor{structure}{fg=accent}
\setbeamercolor{frametitle}{fg=ink, bg=}
\setbeamercolor{normal text}{fg=ink}

% `@` is not a letter in a .tex file, so the emptiness test for the optional
% subtitle has to be wrapped in \makeatletter.
\makeatletter
\setbeamertemplate{frametitle}{%
  \vspace{0.6em}%
  \usebeamerfont{frametitle}\insertframetitle\par
  \ifx\insertframesubtitle\@empty\else
    {\usebeamerfont{framesubtitle}\color{rule}\insertframesubtitle\par}%
  \fi
  \vspace{0.2em}{\color{accent}\rule{\textwidth}{1pt}}\par
}
\makeatother

\setbeamertemplate{footline}{%
  \hbox{\begin{beamercolorbox}[wd=\paperwidth, ht=2.4ex, dp=1.4ex, leftskip=1.6em, rightskip=1.6em]{}%
    {\sffamily\scriptsize\color{rule}\insertshorttitle}\hfill
    {\sffamily\scriptsize\color{rule}\insertframenumber/\inserttotalframenumber}%
  \end{beamercolorbox}}%
}

\setbeamertemplate{itemize item}{\color{accent}\textbullet}
\setbeamertemplate{itemize subitem}{\color{rule}--}
\setbeamertemplate{caption}{\raggedright\sffamily\scriptsize\insertcaption\par}
\setbeamertemplate{blocks}[default]
\setbeamercolor{block title}{fg=accent, bg=}

% One-line takeaway pinned under a figure. Two lines at most, per section 5 of
% docs/quy-uoc-bao-cao.md; the rest is for the presenter to say.
\newcommand{\chotlai}[1]{\vspace{0.3em}{\small #1}}

% Beamer takes the short form in the optional argument; there is no \shorttitle.
\title[Tối ưu hóa hàm mất mát Ridge]{Tối ưu hóa hàm mất mát Ridge}
\subtitle{Dự đoán lãi suất khoản vay trên bộ dữ liệu Lending Club}
\author{Nhóm \dots{} \textbullet{} Lớp Khoa học dữ liệu}
\date{\today}

\begin{document}

% ======================= PHẦN 1. ĐẶT VẤN ĐỀ =======================

\begin{frame}[plain]
  \vspace{1.2em}
  {\usebeamerfont{title}\color{ink}\inserttitle\par}
  \vspace{0.4em}{\color{accent}\rule{0.5\textwidth}{1.2pt}\par}
  \vspace{0.6em}
  {\sffamily\large\color{rule}\insertsubtitle\par}
  \vfill
  {\usebeamerfont{author}\insertauthor\par}
  {\usebeamerfont{date}\color{rule}\insertdate\par}
\end{frame}

\begin{frame}{Dữ liệu và bài toán}{Lending Club, 2007--2018}
  \begin{columns}[T]
    \begin{column}{0.52\textwidth}
      \begin{itemize}
        \item \num{2260668} khoản vay, 151 cột.
        \item Dự đoán lãi suất \texttt{int\_rate}.
        \item Sau làm sạch: $d = 116$ thuộc tính.
        \item Hai quy mô: \num{200000} và \num{1200000}.
      \end{itemize}
    \end{column}
    \begin{column}{0.44\textwidth}
      \begin{itemize}
        \item Loại \texttt{sub\_grade}: một mình nó giải thích $R^2 = 0{,}96$ của lãi suất.
        \item Lending Club tra lãi suất từ bậc tín dụng.
        \item Sau khi loại: $R^2 \approx 0{,}49$.
      \end{itemize}
    \end{column}
  \end{columns}
  \note{Nhấn: rò rỉ không làm sai phép tính tối ưu hóa nào, nhưng làm hỏng phần
    đối chiếu RMSE. Con số 0,96 với 0,49 là thước đo phần rò rỉ.}
\end{frame}

\begin{frame}{Hàm mục tiêu}{Cố định trước khi bắt đầu mọi thí nghiệm}
  \[
    \min_{w \in \R^{d}} \quad f(w) = \frac{1}{2n}\norm{Xw - y}^{2} + \frac{\lambda}{2}\norm{w}^{2}
  \]
  \[
    \nabla f(w) = \tfrac{1}{n} X\trans (Xw - y) + \lambda w,
    \qquad
    \nabla^{2} f(w) = \tfrac{1}{n} X\trans X + \lambda I
  \]

  \begin{itemize}
    \item Lồi mạnh, khả vi vô hạn lần, \textbf{có nghiệm đóng}.
    \item Hessian không phụ thuộc $w$, vì $f$ bậc hai.
    \item $\fstar$ tính được tới sai số máy: $\norm{\nabla f(\wstar)} = 7{,}4\cdot10^{-13}$.
  \end{itemize}
  \note{Vì có f*, mọi hình vẽ f(w_k) - f* chứ không vẽ f. Đây là điều kiện để đo
    sai số tuyệt đối thay vì chỉ so đường cong với nhau.}
\end{frame}

% ======================= PHẦN 2. ĐẶC TRƯNG BÀI TOÁN =======================

\begin{frame}{Ba hằng số quyết định mọi thứ phía sau}
  \[
    L = \lambda_{\max} + \lambda, \qquad \mu = \lambda_{\min} + \lambda, \qquad \kappa = L/\mu
  \]
  \begin{center}\small
    \begin{tabular}{lrrrr}
      \toprule
      Quy mô & $n$ & $L$ & $\mu$ & $\kappa$ \\
      \midrule
      Quét tham số & \num{200000}  & 9{,}1156 & 0{,}034073 & 267{,}5 \\
      Toàn phần    & \num{1200000} & 9{,}0620 & 0{,}034036 & 266{,}2 \\
      \bottomrule
    \end{tabular}
  \end{center}
  \chotlai{Ba hằng số lệch dưới $0{,}6\%$ giữa hai quy mô, nhờ hệ số $\frac{1}{n}$
    trong hàm mục tiêu. Quét lưới trên mẫu nhỏ vì thế hợp lệ.}
  \note{Nếu bỏ 1/n thì L tăng tuyến tính theo n và phải quét lại toàn bộ lưới ở
    mỗi quy mô.}
\end{frame}

\begin{frame}{Chuẩn hóa cột: một can thiệp lên độ khó bài toán}
  \begin{columns}[T]
    \begin{column}{0.54\textwidth}
      \centering
      \resultgraphic[0.68]{scaling-kappa_iters}
    \end{column}
    \begin{column}{0.42\textwidth}\small
      \begin{tabular}{lr}
        \toprule
        Cách mã hóa & $\kappa$ \\
        \midrule
        Thô                & $3{,}87\cdot10^{12}$ \\
        Chỉ trừ trung bình & $1{,}91\cdot10^{12}$ \\
        Chuẩn hóa đầy đủ   & $267{,}5$ \\
        \bottomrule
      \end{tabular}
    \end{column}
  \end{columns}
  \chotlai{Trừ trung bình hạ $\kappa$ hai lần; chia độ lệch chuẩn hạ tiếp bảy tỉ lần.}
  \note{Công nằm ở phép chia chứ không ở phép trừ. Các cột lệch nhau sáu bậc:
    tot_cur_bal cỡ 1e6 bên cạnh dti cỡ hàng chục.}
\end{frame}

% ======================= PHẦN 3. CHỌN ĐỘ DÀI BƯỚC =======================

\begin{frame}{Bước cố định: ngưỡng \texorpdfstring{$2/L$}{2/L} là ngưỡng thật}
  \centering
  \resultgraphic[0.66]{step-fixed_iters}

  \chotlai{$t = 2{,}1/L$ phân kỳ ở vòng 94 vì $\abs{1-tL} = 1{,}1 > 1$;
    $t = 2/L$ dao động không tắt.}
  \note{Đường phân kỳ được giữ trong hình chứ không loại bỏ.}
\end{frame}

\begin{frame}{Bước tối ưu lý thuyết thua bước \texorpdfstring{$1{,}9/L$}{1,9/L}}
  \begin{columns}[T]
    \begin{column}{0.46\textwidth}\small
      \begin{tabular}{lr}
        \toprule
        Bước & Vòng lặp \\
        \midrule
        $t = 1{,}9/L$    & \textbf{331} \\
        $t = 1/L$        & 630 \\
        $t = 2/(L+\mu)$  & 766 \\
        \bottomrule
      \end{tabular}
      \vspace{0.8em}

      \small
      $\abs{1-tL}$: \\
      $1{,}9/L \to 0{,}9000$ \\
      $2/(L+\mu) \to 0{,}9926$
    \end{column}
    \begin{column}{0.5\textwidth}\small
      Sai số ban đầu tại $w_0 = 0$:
      \begin{itemize}
        \item hướng $\lambda > 1$ mang $\mathbf{80{,}5\%}$
        \item hướng $\lambda < 0{,}1$ mang $0{,}3\%$
      \end{itemize}
      \[
        f(w_k) - \fstar = \tfrac{1}{2}\sum_i \lambda_i c_i^2 (1-t\lambda_i)^{2k}
      \]
    \end{column}
  \end{columns}
  \chotlai{Cận lý thuyết là cận cho \emph{trường hợp xấu nhất}; $w_0 = 0$ không
    phải trường hợp đó.}
  \note{Mô hình modal dự đoán 9,80e-7 ở vòng 331 và 9,95e-7 ở vòng 766, khớp
    đúng số quan sát. Kết luận đảo chiều khi cần độ chính xác rất cao.}
\end{frame}

\begin{frame}{Backtracking thắng cả hai trục}
  \begin{center}\small
    \begin{tabular}{lrrr}
      \toprule
      Cấu hình & Vòng lặp & Giây & Lần tính $f$/vòng \\
      \midrule
      Armijo $c=0{,}3$, $\rho=0{,}5$ & \textbf{143} & \textbf{3{,}14} & 4{,}94 \\
      Bước cố định $1{,}9/L$         & 331          & 4{,}75          & 1 \\
      \bottomrule
    \end{tabular}
  \end{center}
  \[
    t \le \frac{2(1-c)\norm{g}^{2}}{g\trans \nabla^{2}f\, g}
  \]
  \chotlai{Armijo so với thương Rayleigh \emph{tại chỗ}, không với $L$ toàn cục,
    nên nhận bước lớn hơn $2/L$ ở hướng phẳng.}
  \note{Ngược dự đoán thông thường là "thắng vòng lặp, thua thời gian".
    Đảo chiều khi mỗi lần tính f đắt ngang một gradient và kappa nhỏ.}
\end{frame}

\begin{frame}{Kết luận về chọn bước}{Áp cho cả bốn thuật toán}
  \begin{itemize}
    \item $\rho$ quan trọng hơn $c$ rất nhiều: sáu cấu hình $\rho = 0{,}5$ đứng
      trên toàn bộ mười hai cấu hình $\rho = 0{,}8$ và $0{,}9$; $c$ đổi 3000 lần
      chỉ làm 167 xuống 143 vòng.
    \item Dò bước bằng tay: trong lưới $10^{-3}$ tới $10$, chỉ $t = 0{,}1$ hội
      tụ. Vùng hội tụ là $(0;\,0{,}219)$.
    \item Tính $L$ và $\mu$ tốn dưới một giây; quét năm bước tốn 195 giây và vẫn
      không trúng bước gần tối ưu.
  \end{itemize}
  \note{Đây là lý do phải tính L thay vì dò. Rho quyết định số lần thử mỗi vòng,
    mà mỗi lần thử tốn một lần tính f.}
\end{frame}

% ======================= PHẦN 4. BA CÁCH ĐÁNH ĐỔI =======================

\begin{frame}{Khung chung cho ba chương tiếp theo}
  \begin{center}
    \large
    thời gian $\;=\;$ số vòng lặp $\;\times\;$ giá mỗi vòng
  \end{center}
  \vspace{0.6em}
  \begin{center}\small
    \begin{tabular}{lll}
      \toprule
      Phương pháp & Đổi đi & Lấy về \\
      \midrule
      Quán tính (AGD) & một vector bộ nhớ & ít vòng lặp hơn 10 lần \\
      Ngẫu nhiên hóa (SGD) & độ chính xác cuối & giá mỗi vòng rẻ $n/B$ lần \\
      Bậc hai (Newton) & giá mỗi vòng đắt $\sim d$ lần & đúng một vòng lặp \\
      \bottomrule
    \end{tabular}
  \end{center}
  \note{Ba phương án nằm ở ba góc khác nhau của cùng một đánh đổi. Chương tổng
    hợp chỉ việc đặt cả bốn lên chung hai trục.}
\end{frame}

\begin{frame}{Quán tính: cùng giá mỗi vòng, ít vòng hơn mười lần}
  \centering
  \resultgraphic[0.68]{momentum-formula_time}

  \chotlai{AGD 286 vòng và $4{,}0$ giây, GD 3000 vòng và $42{,}1$ giây, cùng bước
    $1/L$. $\beta = 0{,}9$ gần đúng ở đây do $\kappa = 267{,}5$ cho
    $\beta^{*} = 0{,}885$.}
  \note{Không nên đọc thành kết luận chung về 0,9. Với kappa = 10 thì
    beta* = 0,52 và 0,9 sẽ sai rõ.}
\end{frame}

\begin{frame}{Khởi động lại: cơ chế bù cho việc không biết \texorpdfstring{$\mu$}{mu}}
  \centering
  \resultgraphic[0.68]{momentum-restart_iters}

  \chotlai{Với $\beta_k = (k-1)/(k+2)$: 143 xuống 89 vòng, đình trệ thành hội tụ.
    Với $\beta$ tính theo $\mu$: trùng nhau từng chữ số.}
  \note{Điều kiện restart không kích hoạt lần nào khi beta đã đặt đúng theo mu.}
\end{frame}

\begin{frame}{Ngẫu nhiên hóa: kích thước lô đổi giá, không đổi sàn}
  \centering
  \resultgraphic[0.70]{batch-size_epochs}

  \chotlai{Mọi $B$ dừng quanh $5\cdot10^{-4}$. $B=32$ tốn \num{250000} vòng và
    $4{,}2$ giây; $B=2048$ tốn 3880 vòng và $2{,}2$ giây.}
  \note{Bước đặt theo 1/L_B bù trừ giữa phương sai gradient và số bước.
    L_max = 200056 so với L = 9,12, gấp 21947 lần, do một dòng ngoại lai.}
\end{frame}

\begin{frame}{Quy tắc giảm bước: kết luận phụ thuộc đúng một tham số}
  \begin{center}\small
    \begin{tabular}{lr}
      \toprule
      Quy tắc, sau 200 epoch & $f - \fstar$ \\
      \midrule
      Bậc thang, giảm nửa mỗi 40 epoch      & $8{,}5\cdot10^{-7}$ \\
      Nghịch đảo, $\gamma = \mu\eta_0$      & $6{,}9\cdot10^{-6}$ \\
      \textbf{Hằng số}                      & $1{,}1\cdot10^{-3}$ \\
      Nghịch đảo, $\gamma = 1/\text{epoch}$ & $1{,}5\cdot10^{-2}$ \\
      Căn bậc hai                           & $4{,}6\cdot10^{-1}$ \\
      \bottomrule
    \end{tabular}
  \end{center}
  \chotlai{Hai dòng nghịch đảo khác nhau đúng một tham số và cho hai kết luận
    ngược nhau. $\gamma$ phải neo vào $\mu$, không vào số vòng mỗi epoch.}
  \note{Ở 40 epoch bước hằng vẫn đang dẫn; chỉ từ khoảng 100 epoch thứ hạng mới
    đảo. Bước hằng còn có phương sai giữa seed lớn hơn 3,7 lần.}
\end{frame}

\begin{frame}{Bậc hai: một vòng lặp, sai số bằng không}
  \[
    w_{1} = w_{0} - (\nabla^{2} f)^{-1}\nabla f(w_{0})
          = \left(\tfrac{1}{n}X\trans X + \lambda I\right)^{-1}\tfrac{1}{n}X\trans y = \wstar
  \]
  \begin{center}\small
    \begin{tabular}{lrl}
      \toprule
      Thành phần & Thời gian & Bậc \\
      \midrule
      Dựng $\frac{1}{n}X\trans X$ & $163{,}2$ ms & $\bigO(nd^{2})$ \\
      Cholesky $116\times116$     & $0{,}2$ ms   & $\bigO(d^{3})$ \\
      Một gradient                & $87{,}3$ ms  & $\bigO(nd)$ \\
      \bottomrule
    \end{tabular}
  \end{center}
  \chotlai{Một bước Newton $\approx 250$ ms, trong đó $\bigO(d^{3})$ chỉ chiếm
    $0{,}2$ ms. Backtracking nhận bước đầy đủ ngay lần thử thứ nhất.}
  \note{Đảo chiều khi d lên vài nghìn: Hessian tăng theo d^2, gradient theo d.}
\end{frame}

% ======================= PHẦN 5. SO SÁNH TỔNG HỢP =======================

\begin{frame}{So sánh tổng hợp: theo số vòng lặp}
  \centering
  \resultgraphic[0.66]{headline_iters}
  \note{Newton 1 vòng, AGD 286, backtracking 561, GD 2022, SGD 3880 nhưng không
    tới đích.}
\end{frame}

\begin{frame}{So sánh tổng hợp: theo thời gian chạy}
  \centering
  \resultgraphic[0.62]{headline_time}

  \chotlai{Thứ hạng \emph{không} đổi giữa hai trục ở bài này, vì giá mỗi vòng
    chênh nhau ít hơn số vòng chênh nhau.}
  \note{Một vòng Newton đắt hơn một vòng GD khoảng ba lần, nhưng GD cần nhiều hơn
    hai nghìn lần số vòng. Thứ hạng sẽ đảo khi d lớn.}
\end{frame}

\begin{frame}{Bảng tổng hợp và đối chiếu lý thuyết}
  \begin{columns}[T]
    \begin{column}{0.5\textwidth}\small
      \begin{tabular}{lrr}
        \toprule
        Cấu hình tốt nhất & Vòng & Giây \\
        \midrule
        Newton            & 1    & 0{,}25 \\
        AGD               & 75   & 1{,}05 \\
        GD backtracking   & 143  & 3{,}14 \\
        GD $t=1{,}9/L$    & 331  & 4{,}75 \\
        SGD               & \multicolumn{2}{c}{không đạt} \\
        \bottomrule
      \end{tabular}
    \end{column}
    \begin{column}{0.46\textwidth}\small
      \begin{tabular}{lrr}
        \toprule
        $\kappa$ giảm & GD & dự đoán \\
        \midrule
        730 $\to$ 268 & 2{,}73 & 2{,}73 \\
        268 $\to$ 90  & 2{,}98 & 2{,}98 \\
        90 $\to$ 10   & 8{,}86 & 8{,}91 \\
        \bottomrule
      \end{tabular}
      \vspace{0.5em}

      GD tỷ lệ với $\kappa$, AGD với $\sqrt{\kappa}$.
    \end{column}
  \end{columns}
  \chotlai{Số vòng lặp không đổi khi $n$ tăng 6 lần; thời gian tăng $5{,}2$ tới
    $6{,}5$ lần, tức tuyến tính theo $n$.}
\end{frame}

% ======================= PHẦN 6. SO SÁNH THƯ VIỆN =======================

\begin{frame}{Quy đổi hàm mục tiêu sang scikit-learn}
  \begin{itemize}
    \item \texttt{Ridge} cực tiểu $\norm{Xw-y}^{2} + \alpha\norm{w}^{2}$
      $\Rightarrow$ $\alpha = \lambda n$.
    \item \texttt{SGDRegressor} cực tiểu đúng $f(w)$ của nhóm
      $\Rightarrow$ $\alpha = \lambda$.
  \end{itemize}
  \vspace{0.4em}
  Hai ước lượng trong cùng thư viện cần hai hằng số khác nhau. Đặt sai
  \textbf{không sinh lỗi}, chỉ lặng lẽ so hai bài toán khác nhau.
  \vspace{0.6em}

  \chotlai{Kiểm chứng bằng số:
    $\norm{\hat{w}-\wstar}/\norm{\wstar} = 2{,}35\cdot10^{-15}$.}
  \note{Phép kiểm này chạy tự động và làm dừng chương trình nếu không đạt.}
\end{frame}

\begin{frame}{Kết quả so với thư viện}
  \begin{center}\small
    \begin{tabular}{lrr}
      \toprule
      Phương pháp, $n = \num{1200000}$ & Giây & $f - \fstar$ \\
      \midrule
      Newton của nhóm            & $0{,}25$  & $0$ \\
      \texttt{Ridge(cholesky)}   & $0{,}273$ & $4{,}2\cdot10^{-30}$ \\
      \texttt{Ridge(lsqr)}       & $2{,}03$  & $1{,}1\cdot10^{-5}$ \\
      \texttt{Ridge(sag)}        & $45{,}4$  & $4{,}0\cdot10^{-7}$ \\
      \texttt{SGDRegressor} mặc định & $3{,}18$ & $6{,}2\cdot10^{20}$ \\
      \bottomrule
    \end{tabular}
  \end{center}
  \chotlai{Hai con số sát nhau ở hai dòng đầu là bằng chứng cả hai cài đúng, vì
    \texttt{Ridge} mặc định cũng là nghiệm đóng qua Cholesky.}
  \note{Mục đích tự cài không phải chạy nhanh hơn thư viện, mà là hiểu cơ chế
    hội tụ và biết chọn tham số. SGDRegressor mặc định không đặt bước theo L.}
\end{frame}

% ======================= PHẦN 7. KẾT LUẬN =======================

\begin{frame}{Hai vai trò của hệ số hiệu chỉnh}
  \begin{columns}[T]
    \begin{column}{0.52\textwidth}
      \centering
      \resultgraphic[0.55]{lambda-kappa_tradeoff}
    \end{column}
    \begin{column}{0.44\textwidth}\small
      \begin{tabular}{rrrr}
        \toprule
        $\lambda$ & $\kappa$ & GD & RMSE \\
        \midrule
        0{,}001 & 2633 & --   & 3{,}4537 \\
        0{,}032 & 268  & 766  & 3{,}4563 \\
        0{,}1   & 90   & 257  & 3{,}4690 \\
        10      & 1{,}9 & 5   & 4{,}4578 \\
        \bottomrule
      \end{tabular}
    \end{column}
  \end{columns}
  \chotlai{Từ $0{,}001$ lên $0{,}1$: RMSE xấu $0{,}44\%$, $\kappa$ giảm 29 lần,
    GD từ không hội tụ xuống 257 vòng.}
  \note{Quy tắc một sai số chuẩn chọn lambda = 0,0316, nằm đúng trong khoảng đó.}
\end{frame}

\begin{frame}{Những điều rút ra}
  \begin{itemize}
    \item \textbf{Chọn tham số.} $\rho$ quan trọng hơn $c$; bước SGD theo
      $1/L_{B}$ chứ không theo $1/L$; tốc độ giảm neo vào $\mu$; $\beta$ tính lại
      theo bước mà line search thực sự chọn.
    \item \textbf{Ba chỗ lệch khỏi lý thuyết}, cùng một lý do: cận lý thuyết là
      cận cho trường hợp xấu nhất, thực nghiệm chạy trên một bài toán, một điểm
      khởi tạo và một ngân sách cụ thể.
    \item \textbf{Phạm vi.} Với $d = 116$ và $n$ cỡ triệu, phương pháp bậc hai
      thắng áp đảo. Phương pháp lặp bậc một chỉ có lợi khi $d$ lớn tới mức không
      dựng nổi $X\trans X$.
  \end{itemize}
  \vspace{0.4em}
  \begin{center}
    {\color{accent}\rule{0.4\textwidth}{1pt}}\par
    \vspace{0.4em}
    {\sffamily Cảm ơn thầy cô và các bạn đã lắng nghe.}
  \end{center}
\end{frame}

\end{document}
